\chapter{引言}
	超声速射流是航空航天推进系统、高超声速风洞实验和天然气超声速分离等工程领域的核心流动形态。在火箭发动机尾喷管出口、超燃冲压发动机燃烧室入口以及高超声速飞行器表面，欠膨胀超声速射流普遍存在，其流场热力学结构——包括温度、压力、密度的空间分布、激波位置与强度、混合层发展等——直接决定设备性能、燃烧效率和气动力特性。准确获取这些流场参数不仅是理解复杂流动现象的基础，更是优化工程设计的先决条件。
	随着高超声速飞行器技术的快速发展，风洞实验中对复杂流场非侵入式诊断的需求日益迫切。在超声速风洞中，为实现流动可视化与定量测量，常需引入示踪物质。二氧化碳（CO$_2$）因其易于获得、安全无毒，且固态干冰可直接升华注入流场作为示踪粒子，成为一种实用选择。然而，在风洞实际运行的低温、高速膨胀条件下，CO$_2$本身会经历从气相到液/固相的非平衡凝结相变。这一相变过程会显著改变局部流动结构、释放潜热并影响温度场，使得示踪粒子的运动轨迹不再完全代表原始流场信息，给流场精准诊断引入不可忽略的误差。
	与此同时，含可凝结组分的高速流动相变问题也是流体力学与热力学交叉领域的研究难点。CO$_2$在超声速膨胀条件下的凝结行为，不仅关乎示踪技术的可靠性，还与碳捕集、天然气提纯等能源工程密切相关。然而，当前对欠膨胀超声速射流中非平衡相变的认知仍存在显著不足：一方面，经典成核理论（CNT）建立在平衡态或准静态假设之上，难以直接描述超声速流场中快速膨胀导致的过饱和与成核过程；另一方面，现有实验研究多采用纹影法等定性手段，缺乏对流场温度、压力等关键热力学参数的同步定量测量，导致相变模型缺少可靠的实验数据支撑。
	从测量技术角度看，传统的流场测量方法各自存在局限。纹影法虽能清晰显示激波等密度梯度结构，但仅能定性观察，无法输出温度、压力等定量参数。粒子图像测速（PIV）依赖外部示踪粒子，而超声速流场对粒子的跟随性要求极高，粒径、密度匹配不当会引入显著的系统误差。接触式探针测量则会破坏流场完整性，且难以应用于高温高速环境。因此，亟需一种非侵入、高灵敏度、快速响应且能直接探测CO$_2$特征的测量方法。
	基于上述需求，\textbf{中红外激光吸收光谱（TDLAS）技术}凭借非侵入式测量、无需添加示踪剂、可直接探测CO$_2$特征吸收谱线的优势，为超声速流场诊断提供了理想方案。该技术基于分子光谱学原理，当激光穿透流场时，CO$_2$分子选择性吸收特定波长的激光，通过测量两条邻近谱线的吸收强度比值，可直接反演温度；结合线宽分析可进一步推导压力和密度，既不破坏流场也不干扰相变过程，是实现流场热力学参数与相变过程同步测量的有效途径。
	\subsection{本文研究内容与核心问题}
	本课题的研究目的为：（1）建立一套基于TDLAS的超声速射流流场温度、压力、密度定量测量方法；（2）揭示欠膨胀射流激波结构对热力学参数分布的调控规律；（3）探究CO$_2$在超声速膨胀条件下的非平衡相变机制及其与流场参数的定量关联。
	基于上述分析，\textbf{本文围绕以下三个核心问题展开}：
	\textbf{核心问题一：如何实现超声速射流流场温度、压力、密度的全分布高精度测量？}欠膨胀射流的本质是CO$_2$工质的绝热膨胀过程，其流场中温度、压力、密度等热力学参数的分布受激波系统的强烈调制，呈现高度非均匀、非平衡的特征。然而，现有测量手段难以在不干扰流场的前提下同步获取这些参数的全场分布，导致流场热力学结构的实验认知长期停留在定性层面。本课题采用中红外激光吸收光谱（TDLAS）结合伺服电机自动扫描和Abel反演的技术路线，通过对轴对称射流进行视线积分光谱扫描，反演获得流场径向参数分布，解决传统测量方法难以同步获取多物理量空间分布的难题。
	\textbf{核心问题二：欠膨胀射流的激波结构对热力学参数分布的影响规律是什么？}欠膨胀射流中马赫盘、桶形激波等典型结构会导致流场参数发生剧烈突变——马赫盘后压力骤升30\%—50\%、温度同步回升，膨胀扇后则进一步降温加速凝结。但激波结构与温度、压力场的定量关联尚不明确，激波/膨胀波系如何调制热力学参数的轴向演化与径向分布，缺乏系统的实验研究。本课题利用纹影技术和TDLAS光谱测量的同步对标，将宏观激波结构与微观温度、压力场建立关联，揭示激波/膨胀波系对温度、压力突变和径向分布的影响机制。
	\textbf{核心问题三：如何建立超声速射流中CO$_2$相变与流场参数的定量关联？}目前主流的CO$_2$相变理论多适用于平衡态环境，而欠膨胀射流中的快速膨胀导致系统处于强非平衡态，现有以准一维膨胀流为代表的模型缺少实验数据验证，不能与温度、压力等热力学参数有效关联。本课题基于准一维相变模型与TDLAS实验测量数据的交叉验证，以经典成核理论（CNT）为基础框架，结合实验测量的温度、压力场数据对CNT关键参数（如表面张力尺度修正、非平衡修正）进行标定与优化，揭示相变发生条件（临界温度、压力、过饱和度）与流场参数的定量关系，为非平衡态相变模型的构建与校准提供实验依据。
	\subsection{全文结构安排}
	本文以“流场测量方法建立→热力学参数分布揭示→相变机制探究”为主线，共分为七章，各章内容安排如下：
	\textbf{第一章 引言。}阐述本课题的研究背景与意义，指出现有超声速流场诊断技术在温度、压力、密度同步定量测量方面的不足，以及CO$_2$非平衡相变研究中理论与实验脱节的问题。在此基础上，明确提出本文的三个核心研究问题：如何实现超声速射流流场热力学参数的全分布高精度测量？欠膨胀射流的激波结构对热力学参数分布的影响规律是什么？如何建立超声速射流中CO$_2$相变与流场参数的定量关联？最后给出全文结构导航。
	\textbf{第二章 文献综述。}系统综述欠膨胀射流流场结构与测量、激光诊断技术发展、CO$_2$相变机制三个方面的研究进展。在欠膨胀射流部分，梳理从Shapiro的开创性理论到Crist、Addy等人的经典实验，涵盖马赫盘位置与直径的定量规律、激波胞结构特征、微尺度效应及真实气体模型修正等核心成果。在激光诊断技术部分，对比分析PLIF、PIV、纹影法、分子束光谱及中红外激光光谱测量技术的适用范围与局限性，重点阐述TDLAS技术在非侵入式温度、压力、密度同步测量方面的优势及关键数据处理方法（Abel反演、双线强度比法、Voigt线型拟合）。在相变机制部分，从Gibbs界面热力学出发，追溯经典成核理论（CNT）从Volmer-Weber到Becker-Döring的理论发展脉络，介绍Harding的数值计算程序及Wu、Li等人在欠膨胀射流CO$_2$团簇可视化方面的近期进展，指出现有研究在非平衡相变与流场参数定量关联方面的不足，明确本课题的研究定位。
	\textbf{第三章 Fluent仿真与硬件系统。}本章分为两大部分。第一部分介绍Fluent流场仿真方法，包括二维轴对称计算模型、结构化网格划分、边界条件设置（压力入口0.5~MPa、压力出口0.1~MPa）、SST $k$-$\omega$湍流模型及耦合求解器配置，给出温度场、压力场、速度场的云图结果，提取沿轴线和马赫盘附近截面的参数径向分布，为后续实验提供工况依据和验证真值。第二部分详细介绍实验硬件系统的设计与搭建，包括射流发生系统（高压气源、混合腔体、小孔喷嘴）、TDLAS双光路系统（2.7~$\mu$m QCL激光器、HgCdTe探测器、光路布局与工作参数）、反射式Z型纹影系统、伺服电机位移扫描系统及Red Pitaya数据采集系统，给出各子系统的设备型号、关键参数和调试结果。
	\textbf{第四章 光谱仿真与纹影诊断。}首先基于HITRAN数据库进行line-by-line光谱仿真，覆盖200—300~K、0.03—0.5~MPa的工况范围，通过对比2.7~$\mu$m、2~$\mu$m、4.2~$\mu$m三个波段的吸收线强、温度灵敏度和H$_2$O干扰程度，确定最佳实验波段和测温谱线对（$\nu_1=3694.3$~cm$^{-1}$，$\nu_2=3696.2$~cm$^{-1}$，$\Delta E''\approx660$~cm$^{-1}$）。其次介绍FSR标准具波长标定方法，建立采样点-波数映射关系。第三，通过高斯解析场加噪测试对比两点法、三点法和洋葱剥皮法的精度与抗噪性能，给出误差统计结果，确定三点法为实测数据核心反演算法。第四，给出纹影系统的诊断结果，包括射流试验台实物照片、不同上游压力下的纹影图像，马赫盘位置与Fluent仿真对比验证。
	\textbf{第五章 层析扫描实验与温度分布。}展示纯CO$_2$欠膨胀射流（0.5~MPa，1~mm喷嘴）的层析扫描实验结果。首先给出3轴向平面（$Z=10$、20、30~mm）$\times$6径向点（$R=0$—5~mm）的扫描方案和实验条件。其次展示$Z=20$~mm平面不同径向位置的原始吸收光谱，观察中心最强、径向递减的趋势。第三，通过三点法Abel反演得到径向吸收率分布，识别射流半宽约1.6~mm、边界约3.8~mm。第四，采用双线强度比法反演温度，给出径向温度分布表，中心268~K、径向递增至305~K。第五，将实验结果与Fluent仿真和Wu等（2021）文献结果进行对比验证。第六，进行误差分析，给出综合不确定度约±6~K（95\%置信区间）。
	\textbf{第六章 结果分析与相变模型。}对温度反演结果进行深入分析，讨论各误差来源的贡献及改进方向。基于准一维相变模型（Wu等2021代码），以实测温压场为输入，预测CO$_2$团簇成核位置和生长特征。给出模型输入参数（进口总压0.5~MPa、总温300~K、纯CO$_2$、直孔1~mm）、计算流程、轴线温度/压力演化曲线、凝结起始位置（$z\approx0.8$~mm）、临界半径（$r^*\approx1.2$~nm）、成核速率峰值位置（$z\approx1.5$~mm，$J\approx10^{20}$~m$^{-3}$s$^{-1}$）及马赫盘后团簇蒸发趋势。将预测结果与Wu等（2021）实验（$z/D\approx0.1-0.4$）对比，分析偏差原因（喷嘴几何差异、过冷度判定阈值不同），指出后续校准方向。
	\textbf{第七章 结论与展望。}总结全文的主要研究成果，包括：（1）TDLAS层析温度测量方法的建立与验证；（2）欠膨胀射流温度场径向分布规律的初步揭示；（3）实验系统硬件改进的关键成果（聚焦0.3~mm、噪声降80\%、干扰降96\%）；（4）CNT模型基于实验数据的初步预测。在此基础上提炼本课题的创新点，指出研究存在的局限性（如尚未直接探测到相变信号、扫描点数有限、压力反演精度待提高、相变模型未收敛），并对下一步工作提出展望：扩大扫描范围以覆盖相变关键区域、优化FSR波长标定算法、引入FPGA硬件平均提升信噪比、完善相变模型与实验数据的定量对比分析。
% ==================== 第二章 文献综述 ====================

